\subsubsection{Diferencia de umbral y composición de coronas}
\label{mg07:umbral}

La ampliación de una célula discreta conserva todos los estratos que
aparecen al aumentar su tamaño en una unidad. La diferencia de umbral
expresa esa operación sobre el lector de cardinales y permite seguir
su comportamiento bajo suma y producto antes de introducir una
realización métrica.

\begin{definition}[Diferencia de umbral]
Sea \(R\) un anillo conmutativo con unidad y sea
\(f:\mathbb Z_{\geq0}\to R\). Definimos
\begin{equation}
 (D_\#f)(n)=f(n+1)-f(n).
 \label{mg07:definicion}
\end{equation}
Para un lector entero de células finitas anidadas, esta diferencia
registra el cardinal de la corona incorporada en el paso \(n\to n+1\).
\end{definition}

\begin{proposition}[Aditividad y ley exacta de producto]
Para \(f,g:\mathbb Z_{\geq0}\to R\), se cumplen
\begin{align}
 D_\#(f+g)(n)&=D_\#f(n)+D_\#g(n),
 \label{mg07:aditividad}\\
 D_\#(fg)(n)&=f(n+1)D_\#g(n)+g(n)D_\#f(n)
 \label{mg07:producto-desplazado}\\
 &=f(n)D_\#g(n)+g(n)D_\#f(n)
   +D_\#f(n)D_\#g(n).
 \label{mg07:producto}
\end{align}
\end{proposition}

\begin{proof}
La diferencia de la suma se separa término a término. Para el producto,
introducir y sustraer \(f(n+1)g(n)\) da
\[
 \begin{split}
 f(n+1)g(n+1)-f(n)g(n)
 &=f(n+1)\bigl(g(n+1)-g(n)\bigr)\\
 &\quad+g(n)\bigl(f(n+1)-f(n)\bigr).
 \end{split}
\]
Esto prueba \eqref{mg07:producto-desplazado}. La sustitución
\(f(n+1)=f(n)+D_\#f(n)\) produce
\eqref{mg07:producto} con su término bilineal completo.
\end{proof}

Si \(A_n\subseteq A_{n+1}\) y \(B_n\subseteq B_{n+1}\) son
conjuntos finitos, la corona del producto se descompone en tres partes
disjuntas:
\[
 \begin{split}
 (A_{n+1}\times B_{n+1})\setminus(A_n\times B_n)
 &=\bigl(A_n\times(B_{n+1}\setminus B_n)\bigr)\\
 &\quad\sqcup\bigl((A_{n+1}\setminus A_n)\times B_n\bigr)\\
 &\quad\sqcup\bigl((A_{n+1}\setminus A_n)
                 \times(B_{n+1}\setminus B_n)\bigr).
 \end{split}
\]
Tomar cardinales da exactamente \eqref{mg07:producto}. El término
bilineal cuenta las posiciones nuevas en ambos factores; su residencia
es la intersección de las dos ampliaciones, con la incidencia conservada.

\begin{proposition}[Ley binomial de la corona dimensional]
Para \(d\geq1\), sea \(f_d(n)=n^d\). Entonces
\begin{equation}
 D_\#f_d(n)=(n+1)^d-n^d
 =\sum_{k=0}^{d-1}\binom dk n^k
 =\sum_{r=1}^{d}\binom dr n^{d-r}.
 \label{mg07:binomio}
\end{equation}
\end{proposition}

\begin{proof}
La expansión binomial de \((n+1)^d\) contiene \(n^d\) y los
\(d\) términos restantes de la primera suma. La sustitución
\(r=d-k\) da la segunda. En el soporte
\(\{1,\ldots,n+1\}^d\), cada posición añadida posee un conjunto
no vacío de coordenadas iguales a \(n+1\). Si ese conjunto tiene
\(r\) elementos, hay \(\binom dr\) elecciones de posiciones y
\(n^{d-r}\) elecciones de las coordenadas anteriores. Estos
estratos son disjuntos y cubren la corona; por tanto, también prueban
la fórmula mediante incidencia finita.
\end{proof}

Para \(d=3\), los tres estratos tienen cardinales \(3n^2\),
\(3n\) y \(1\). En el paso nonádico--decimal, \(n=9\), la
operación da \(243+27+1=271\); la exclusión del vértice final deja
\(270\) posiciones. La completación cúbica que sigue explicita la
unión de esta corona con las \(729\) posiciones anteriores. La ley
de umbral conserva así tanto el total como la distribución por número
de coordenadas nuevas. Las realizaciones reticulares posteriores
transportan sus incidencias mediante los mapas específicos de cada
construcción.
